% Lösungen Vektoren und Rechenoperationen (zusammenbau v0.4, Heft)
\einheitenkopf{Vektoren und Rechenoperationen}

\erg{19}{a) Zahl \quad b) $\overrightarrow{AB} = (3 | 1 | 2)$, $\overrightarrow{BA} = (-3 | -1 | -2)$, $2 \cdot \overrightarrow{AB} = (6 | 2 | 4)$ \quad c) $1 + a^2 + 4 = 9$, $a^2 = 4$, wegen $a < 0$ ist $a = -2$}
\erg{20}{a) $25 + 144 + (a - 3)^2 = 169$, also $(a - 3)^2 = 0$ und $a = 3$ – nur eine Lösung \quad b) Ja: $(600 | 0 | 0)$ hat den Betrag 600, $(200 | 200 | 200)$ den Betrag $\sqrt{120\,000} \approx 346{,}4$ – gleiche Summe, verschiedene Längen. \quad c) Nein: $(90 | 0 | 0)$ hat den Betrag 90, $(30 | 30 | 30)$ den Betrag $\sqrt{2700} \approx 51{,}96$.}
\erg{21}{a) Ecke ($F$) \quad b) $P(3 | 4 | 0)$, Mitte der Kante $DC$ \quad c) $r = 0$, $s = 1$, $t = 1$, also $\overrightarrow{M_1M_2} = \vec v + \vec w$}

21b)

\begin{ksys3}[x1max=7] \rquader{0,0,0}{6}{4}{3} \rpunkt*{0}{0}{0}{A} \rpunkt*{6}{0}{0}{B} \rpunkt*{6}{4}{0}{C} \rpunkt*{0}{4}{0}{D} \rpunkt*{0}{0}{3}{E} \rpunkt*{6}{0}{3}{F} \rpunkt*{6}{4}{3}{G} \rpunkt*{0}{4}{3}{H} \rpunkt*{3}{4}{0}{P} \end{ksys3}

\erg{22}{a) $P(3 | 2 | 3)$, Mitte der Deckfläche $EFGH$ \quad b) $P(6 | 2 | 3)$, Mitte der Kante $FG$ \quad c) Mitte der Seitenfläche $ADHE$: $\overrightarrow{AG} - \overrightarrow{AB} = \overrightarrow{BG}$ ist die Flächendiagonale, ihre Hälfte führt von $A$ aus in die Flächenmitte – nicht auf eine Kante. \quad d) $D(1 | 3 | 0)$, denn $\overrightarrow{OD} = \overrightarrow{OA} + \overrightarrow{BC}$ \quad e) $H'(-2 | 2 | 2)$, denn $S(4 | 2 | 3)$ ist die Mitte von $AG$ und jeder Punkt wird um $-\overrightarrow{OS}$ verschoben}

22a)

\begin{ksys3}[x1max=7] \rquader{0,0,0}{6}{4}{3} \rpunkt*{0}{0}{0}{A} \rpunkt*{6}{0}{0}{B} \rpunkt*{6}{4}{0}{C} \rpunkt*{0}{4}{0}{D} \rpunkt*{0}{0}{3}{E} \rpunkt*{6}{0}{3}{F} \rpunkt*{6}{4}{3}{G} \rpunkt*{0}{4}{3}{H} \rpunkt*{3}{2}{3}{P} \end{ksys3}


22b)

\begin{ksys3}[x1max=7] \rquader{0,0,0}{6}{4}{3} \rpunkt*{0}{0}{0}{A} \rpunkt*{6}{0}{0}{B} \rpunkt*{6}{4}{0}{C} \rpunkt*{0}{4}{0}{D} \rpunkt*{0}{0}{3}{E} \rpunkt*{6}{0}{3}{F} \rpunkt*{6}{4}{3}{G} \rpunkt*{0}{4}{3}{H} \rpunkt*{6}{2}{3}{P} \end{ksys3}

\erg{23}{a) $\overrightarrow{OF'} = \overrightarrow{OB} + \frac{2}{5} \cdot (4 | 3 | 0)$: $(4 | 3 | 0)$ steht senkrecht auf $\overrightarrow{AB}$, liegt in der $x_1x_2$-Ebene, zeigt in positive $x_1$-Richtung und hat die Länge 5; der Faktor bringt ihn auf die Länge $|\overrightarrow{BF}|$. \quad b) $\overrightarrow{OF'} = \overrightarrow{OB} + \frac{4}{5} \cdot (3 | 4 | 0)$: $(3 | 4 | 0)$ steht senkrecht auf $\overrightarrow{AB}$, liegt in der $x_1x_2$-Ebene, zeigt in positive $x_1$-Richtung und hat die Länge 5; der Faktor bringt ihn auf die Länge $|\overrightarrow{BF}|$.}
\erg{24}{a) Vektor \quad b) $(12 | 30 | 18)$; die zweite Komponente zählt die verkauften Cappuccinos. \quad c) $42 + 36 + 30 = 108$; der Einkauf kostet 108 €.}
\erg{25}{a) Die Gesamtkosten des Einkaufs in Euro – eine Zahl, kein Vektor. \quad b) $\vec m = t \cdot (5 | 2 | 1)$, $t \cdot (20 + 12 + 8) = 168$, $t = 4{,}2$; $w = 21$, $g = 8{,}4$, $r = 4{,}2$ \quad c) $\vec m = t \cdot (6 | 1 | 2)$, $t \cdot (18 + 12 + 10) = 280$, $t = 7$; $h = 42$, $n = 7$, $c = 14$}
